\chapter[Kurztitel]{Langtitel: Welcher nur bei der Kapitel�berschrift erscheint}
\writer{Sebastian Hoffmann}

\section{Einleitung}

Die einf�hrende Literatur f�r dieses Seminarthema stammt aus \emph{Introduction to Algorithms} \cite[Kapitel 1ff.]{Cormen} von \citeauthor{Cormen} aus dem Jahre \citeyear{Cormen}. Weitere Informationen �ber das Zitieren mit dem Natbib Paket finden sich unter: 
\begin{center}
\texttt{http://merkel.zoneo.net/Latex/natbib.php}
\end{center}
Das Standardwerk \cite{Cormen} ist folgenderma�en in der sample.\texttt{bib}-Datei gespeichert und �ber das Schl�sselwort \texttt{Cormen} verf�gbar:

\begin{verbatim}
@book{Cormen,
 author = {Cormen, Thomas H. and
           Stein, Clifford and
           Rivest, Ronald L. and
           Leiserson, Charles E.},
 title = {Introduction to Algorithms},
 year = {2001},
 isbn = {0070131511},
 edition = {2nd},
 publisher = {McGraw-Hill Higher Education},
}
\end{verbatim}


\section{Hauptteil}

In diesem Teil schauen wir uns ganz kurz ein paar \LaTeX Befehle an, die f�r mathematische Formulierungen notwendig sind. Da es dazu aber zahlreiche Dokumente im Internet gibt, gibt es hier eine altbekannte Formel $a^2+b^2=c^2$ lediglich in ein paar verschiedenen Darstellungen
$$a^2+b^2=c^2$$
welche man nat�rlich auch nummerieren und mit einem Schl�sselwort versehen kann
\begin{equation}
\label{eq:pythagoras}
a^2+b^2=c^2
\end{equation}
um darauf wieder verweisen zu k�nnen -- siehe \eqref{eq:pythagoras}. Man kann das nat�rlich auch als Skalarprodukt darstellen
$$\begin{pmatrix} a & b \end{pmatrix} \cdot \begin{pmatrix} a \\ b \end{pmatrix} = c^2$$
oder mit einer Tabelle schemenhaftes ein rechtwinkliges Dreieck andeuten
\begin{center}
\begin{tabular}{clcc}
    & $\vdots\ddots$ \\
$a$ & $\vdots$ & $\ddots$ & $c$\\
    & $\vdots\hdots$ & $\hdots$ & $\ddots$\hspace{-5mm}$\hdots$ \\
    &          & $b$ \\
\end{tabular}
\end{center}


\section{Schlussteil}

F�r die Kompilierung mit \texttt{latex} -- nicht \texttt{pdflatex} -- nutzen wir das Einf�gen von Bildern im \texttt{eps}-Format:

\begin{verbatim}
\begin{figure}[htb]
\centering\includegraphics[keepaspectratio,width=4cm]{Logo.eps}
\caption{... (siehe Abbildung \ref{fig:logo})...
\label{fig:logo}
\end{figure}
\end{verbatim}

\begin{figure}[htb]
\centering\includegraphics[keepaspectratio,width=4cm]{Logo.eps}
\caption{Das ist das Logo des Instituts f�r Informatik der JGU, worauf wir auch mit einem einfachen Aufruf verweisen k�nnen (siehe Abbildung \ref{fig:logo}). Die Abbildung wird flie�end in den Text eingebunden, und zwar nach den Priorit�ten \textbf{h}ere, \textbf{t}op, \textbf{b}ottom.}
\label{fig:logo}
\end{figure}

Dieser Quellcodeschnipsel steht direkt darunter, aber es kann passieren, dass \LaTeX das Bild hier nicht sinnvoll platzieren kann und es deshalb versucht an den Anfang (dann Ende) der n�chsten Seite zu stellen.

\section{Anhang}

Viel Spa� beim Experimentieren!
